\section{Symbolism and Metaphysics}

\begin{quotex}
Natural phenomena in general, and especially astronomical phenomena, are never looked upon by the traditional doctrines other than as a simple means of expression, whereby they symbolize truths of a higher order.
\flright{\textsc{Rene Guenon}, \emph{Man and his Becoming according to the Vedanta}}
\end{quotex}


\paragraph{Inner States}


It is reasonably straight-forward to describe the sensual, or physical, world. For example, I can describe the banana tree I see outside my window. I can specify its height, the number, size, and shape of the fronds and how they are cut by the wind. I can write about the bunches of fruit and whether they appear ripe.

However, when it comes to inner states, such a description is not so straight forward. This is exacerbated by the fact that very few men have been trained to observe their inner states in any detail, other than as the most primitive emotions, desires, and sensations. Thus, for the initiated, recourse is made to natural phenomena to express subtle inner states. This is furthermore justified since the natural or physical world is the reflection of the supernatural. Guenon explains:

\begin{quotex}
If natural phenomena symbolize such truths of a higher order , it is because their laws are fundamentally nothing but the expression of these very truths in a particular domain, a sort of translation of the corresponding principles, naturally adapted to the special conditions of the corporeal and human state.
\end{quotex}


In this regard, we can do no better than to repeat Guenon's emphasis, ``\emph{It is absolutely essential for the understanding of these matters}.''

\paragraph{Lunar and Solar}


In alchemy, chemical reactions are used to symbolize inner states and transformations. This seems to be understood by most readers of alchemical texts, except for those who seem to be more impressed by extraordinary physical phenomena. \emph{A fortiori}, astronomical symbolism, as we pointed out previously\footnote{\url{https://www.meditationsonthetarot.com/the-cosmic-hierarchy}}, must be understood the same way. As Guenon writes:

\begin{quotex}
It must be clearly understood, that when mention is made, for example, of the Spheres of the Sun and of the Moon, it is never the sun and the moon as visible stars, belonging purely to the corporeal realm, that are referred to, but rather the universal principles which these stars represent after their own fashion in the sensible world ... Indeed the different Worlds, planetary Spheres and elementary Realms which are symbolically described as so many regions (only symbolically however, since the being who journeys through them is no longer subject to space), \textbf{are in reality just different states}.
\end{quotex}


{}


The terms Solar and Lunar themselves have been badly misunderstand, largely in part to \textbf{Julius Evola}'s adoption of the terms from the anthropologist, \textbf{J J Bachofen}. This has led to the opinion that they represent a choice, rather than inner states. Guenon writes that the heart corresponds to the Sun. Of course, metaphysically, the Heart is the seat of intelligence, not of feelings as it is today. The Moon, then, corresponds to the brain, the center of discursive thought which is the reflection of the Intellect. To try to be one or the other is to be imbalanced, not to say impossible; their proper relationship has been described at The Word Is Made Flesh\footnote{\url{https://www.meditationsonthetarot.com/the-word-is-made-flesh}}.

Just as we should not confuse the metaphysical Sun and Moon with their astronomical analogues, neither should we confound them with merely human genders. The temptation is to use Solar and Lunar as mutually exclusive labels, which is tantamount to saying that the Heart can live without the Brain.

\paragraph{The Sky God}


It is very difficult for most men to think abstractly, beyond sensual representations of objects in space. Guenon writes about ``the error of those who imagine they believe that the doctrines in question are only intended to describe and explain phenomena.'' This error is most apparent in men who think they are ``sophisticated'' because they refuse to believe in a ``sky god''. The error is compounded when they replace it by a ``life force'' or an ``immanent god'', and the like. The real ``solution'' is to understand what the symbols mean. Guenon, for example lists these:

\begin{table}[h]\centering\small
\begin{tabular}{cc}\toprule
\textbf{Symbol} & \textbf{Meaning} \\\midrule Earth & Gross manifestation \\ Atmosphere & Subtle manifestation \\ Sky & Formless manifestation \\ & The unconditioned beyond all manifestation \\\bottomrule
\end{tabular}
\end{table}

So a reference to a god in the sky, means that God is beyond all manifestation. The alleged dichotomy between a transcendent and immanent god is simply meaningless. The Atma, or Self, dwells in the Heart, beyond all manifestation. Yet, the Atma is our very being. While Atma is Brahma, beyond all manifestation, as Purusha, it indwells man as his vital centre. This centre is beyond both the light of the Sun and Moon, as it is self-luminous. This is hardly abstract language, but rather a call to recognize it in one's own consciousness. Beyond every change in consciousness, what is it that is always there and never changes?

\paragraph{Tantra}


In his \textit{Foundations of Tibetan Mysticism}, \textbf{Lama Anagrika Govinda} explains the anthropomorphic symbolism of the Tantras:

\begin{quotex}
The abstractness of philosophical concepts and conclusions requires to be constantly corrected by direct experience, by the practice of meditation and the contingencies of daily life. The anthropomorphic element in the Vajrayana is therefore not born from a lack of intellectual understanding (as in the case of primitive man), but, on the contrary, from the conscious desire to penetrate from a merely intellectual and theoretical attitude to the direct awareness of reality. This cannot be achieved through building up convictions, ideals, and aims based on reasoning, but only through conscious penetration of those layers of our mind which cannot be reached or influenced by logical arguments and discursive thought.
\end{quotex}


If logical arguments and discursive thought do not suffice to describe the direct experience of these inner states, then we must resort to symbols. The Lama continues:

\begin{quotex}
Such penetration and transformation is only possible thorough the compelling power of inner vision, whose primordial images or archetypes are the formative principles of our mind. Like seeds they sink into the fertile soil of our subconsciousness in order to germinate, to grow and to unfold their potentialities.  Such visions are not hallucinations, because their reality is that of the human psyche. They are symbols, in which the highest knowledge and the noblest endeavor of the human mind are embodied. Their visualization is the creative process of spiritual projection, through which inner experience is translated into visible form.
\end{quotex}


In particular, this applies to the image of yabyum\footnote{\url{https://www.gornahoor.net/?p=3231}}, which is mistaken as a technique used to achieve enlightenment. Quite the contrary, this symbolizes the union of the male and female principles in consciousness. Lama Govinda writes:

\begin{quotex}
If the passive, all-embracing female principle, from which everything proceeds and into which everything recedes, is united with dynamic male principle of active universal love and compassion, then perfect Buddhahood is attained.
\end{quotex}


This union of the heart and the head can only be represented symbolically, at least for those who have yet to experience that realization. The Lama continues:

\begin{quotex}
The process of Enlightenment is therefore represented by the most obvious, the most human and at the same time, the most universal symbol imaginable: the union of male and female in the ecstasy of love, in which the active element is represented as a male and the passive by a female figure.
\end{quotex}


Those who insist on some sort of Sex Magick based on Tantric or alchemical writings have missed the point. Lama Govinda emphasizes this point:

\begin{quotex}
We must not forget that the figureal representations of these symbols are not looked upon as portraying human beings, but as embodying the experiences and visions of meditation. In such a state, however, there is nothing more that could be called sexual; there is only the super-individual polarity of all life, which rules all mental and physical activities, and which is transcended only in the ultimate state of integration, in the realization of \emph{sunyata}.
\end{quotex}



\begin{center}* * *\end{center}

\begin{footnotesize}
\begin{sffamily}
\texttt{J on 2012-02-13 at 21:29 said:}

``Those who insist on some sort of Sex Magick based on Tantric or alchemical writings have missed the point.''

Who are these people? Seems like a straw man to me. Beyond that, is it not the point of sex magic to go beyond the material, to see the act and those involved as symbols. not unlike Arjuna sees the enemy on the battlefield? I fail to see why this should be an impossible or necessarily incorrect path.

\hfill

\texttt{Cologero on 2012-02-13 at 21:52 said:}

Mr J,

A ``straw man'' is an argument that misrepresents a position. First of all, we eschew arguments. But for the sake of argument, we have not misrepresented anything, as there are indeed those who practice Sex Magick.

It is obvious to anyone who has actually read and understood the post that the sex act is a symbol. We also pointed out what it symbolizes. Do you disagree with Lama Govinda in that regard?

Since the yabyum symbolizes the union of opposites (yan and yang), sex magick is not at all necessary to achieve that union, which is the point being missed.

Now, if you think that sex magick is nevertheless useful in some cases, then please tell us why or share your experiences. It is certainly not sufficient, or else the entire human race would be enlightened by now.

But to return to the alchemical analogy, chemists may deal with mercury and sulfur all day long and never see them as a symbol. Something more must be added, don't you think? And once you understand the sex act is a symbol for something else, how many more times do you have to repeat it to be convinced?

\hfill

\texttt{J on 2012-02-13 at 22:50 said:}

First let me make my intentions plain, and insist against any inclination on your part to see me as a mere disturber. I'm genuinely confused as to what your point is and actually do want to understand you, as I've enjoyed your writings thus far.

Unfortunately, your comment has only confused me more, the greater for its tone. Perhaps rather than responding with snide arrogance, one could fathom the possibility that clarity to you is not necessarily an objective matter for the rest of us, Mr. Cologero, despite the obvious care taken in checking these posts over for mistakes, such as the moon corresponding to the moon. Perhaps I'm too ignorant to grasp the meaning of that one.

It would seem to me that by ``Sex Magic'' you mean two different things.

``It is certainly not sufficient, or else the entire human race would be enlightened by now.''

This is simply unfair. Clearly the entire human race has sex, but do they all claim to be practicing sex magic?

``sex magick is not at all necessary to achieve that union''

And so you would proscribe it for not being the only way? There are differing natures, I'm sure you're aware, and ways for each of them to reach the divine, not all of them intellectual in nature, requiring such studies and practices as you demand.

It would seem that you want to move from unnecessary to unorthodox or incorrect. If I've mistaken your motive, please set me aright, as ``Those who insist on some sort of Sex Magick based on Tantric or alchemical writings have missed the point'' seems to me to be in this spirit, and also to miss the point in another direction. It seems to me to assume an ignorance of the principles involved in those who might take an active involvement in sexual activity as a means of realization. This is clearly beyond your capacity to know.

As for ``how many times must you repeat it'', this seems quite beside the point. Would you have the enlightened commit suicide for having finished the task of awakening? Should one stop yoga, meditation, etc. once satisfied that he's achieved perfect knowledge?

As for personal experiences with ``sex magic'', I have made no attempt at such a thing, so regretfully have nothing to offer in that regard, though I'm sure such findings would have been treated with utmost respect.

\hfill

\texttt{Cologero on 2012-02-13 at 23:55 said:}

Thanks for pointing out the typo.

The point of the post is quite clear and is independent of the virtues or vices of the practice of Sex Magick.

We have previously offered some impressions of
\href{http://www.gornahoor.net/?p=51}{Sex Magick}.


You did bring up an interesting point even if by accident: To be enlightened is to have already committed suicide.

\hfill

\texttt{Charlotte Cowell on 2012-02-14 at 18:47 said:}

We're back to the reunification of soulmates, the miracle greater than the parting of the Red Sea. I could write a book about this topic but will spare you!

\hfill

\texttt{Charlotte Cowell on 2012-02-14 at 18:49 said:}

The Zohar and its professional commentators are probably better

\url{http://alchemical-weddings.com/alchemical-weddings/reunification-of-soul-mates}


When their time to be married arrives, the Creator who knows these spirits and souls joins them as they were before they came into the world.. When they are joined together, they become one body and one soul, right and left in proper unison. As it is written, ``There is nothing new under the sun,'' because this is nothing new but a return to how they were before coming down to this world.

You might say, ``But we have learned that a man obtains a soul mate according to his deeds and his behaviour.'' It is assuredly so! If he is meritorious and his ways are correct, then he deserves his on soul mate, to join her as they were joined when they left the Creator, and before bonding with a physical body.''

Nothing that takes place in the physical dimension is more celebrated by Kabbalah than the joining of soul mates, and toward that end almost any hardship or sacrifice is permissible.

Cx

\end{sffamily}
\end{footnotesize}

